\section{Stability of the actual return records}
\label{sec:return-stability}
We compare the first returns themselves, including their changing domains.
The comparison is made on the common collision space
$M=\T\times(-1,1)$ with the probability $\nu$ in \eqref{eq:flux}.
It is not a comparison between independently chosen inducing schemes.
The squares defining $Y_R$ move analytically with $R$, and $U$ is fixed.
Consequently $Y_R^*$ has constant measure $c_*$, and its boundary is a
finite union of smooth arcs, continuously varying with $R$.

We use only ergodicity, not a parameter-uniform mixing rate, from the
finite-horizon Sinai theory. In the present family the obstacles are
smooth, strictly dispersing and disjoint, and the horizon is finite.
The flow-mixing theorem \cite[Corollary 1.3 and correction note]{BDL}
therefore implies ergodicity of the collision map: the suspension of a
nontrivial invariant collision set would be a nontrivial invariant flow
set. Thus every $Y_R^*$ has full saturation. This explicitly supplies the
full-tower hypothesis in the Kac identity used above.

Write $r=r_R^*$, $F=F_R^*$ and $\tau^*=\tau_R^*$ in sums indexed by
induced returns. For $n\ge1$ define the actual cumulative record
\[
 J_{n,R}=\left(\sum_{j<n}\kappa_R^*\circ F^j,
                   \sum_{j<n}r_R^*\circ F^j,
                   \sum_{j<n}\tau_R^*\circ F^j\right)
       =(K_{n,R},N_{n,R},T_{n,R})
\]
on $Y_R^*$. Its values lie in $\Z^3\times\R$; in particular the last
coordinate is one cumulative time, not $n$ independent roof coordinates.
A tilde denotes extension by zero to $M$. Let $p_{n,R}$ denote the raw
mixed density of $J_{n,R}$ under $\nu_R^*$, whose existence is included
in Theorem~\ref{thm:return-stability} below. Norms of densities use
counting measure times Lebesgue measure, denoted by $\mathfrak m$.

\begin{lemma}[Almost-everywhere stability of a finite return record]
\label{lem:finite-return-stability}
Fix $R_0\in I$ and a finite $n$. Outside a $\nu$-null subset of $M$,
$\widetilde J_{n,R}(x)\to\widetilde J_{n,R_0}(x)$ as $R\to R_0$.
For $x\in Y_{R_0}^*$ outside that null set, the physical word, every
intermediate return decision, $K_{n,R}(x)$ and $N_{n,R}(x)$ are constant
for all sufficiently close $R$. On a neighborhood of such a point the
cumulative time is the physical analytic length of that fixed word.
The neighborhood and parameter tolerance may depend on $x$ and $n$.
\end{lemma}
\begin{proof}
Remove the boundary of $Y_{R_0}^*$, all singular initial states and all
preimages under $T_{R_0}$ of that boundary. These sets are null: the
boundary has zero area, the map preserves $\nu$, and the billiard
singularities and their iterates have zero flux measure. Remove also
the null set on which a required return is infinite. For a remaining
$x\in Y_{R_0}^*$, the first $n$ returns contain finitely many regular
physical collisions. Every tested state lies strictly inside or strictly
outside the return set. The finite-itinerary analyticity following
\eqref{eq:entry}, and the continuous motion of the return boundary,
preserve all these decisions for nearby $(R,x)$. They also preserve
the selected entry roots and lattice labels. The final physical count
and displacement are therefore unchanged, while the sum of flight
lengths varies analytically. A remaining point outside $Y_{R_0}^*$
stays outside $Y_R^*$ for nearby $R$, and its zero extension stays zero.
These statements prove the asserted convergence. No common neighborhood
for infinitely many words is used.
\end{proof}

\begin{theorem}[Finite-record stability in density and first moment]
\label{thm:return-stability}
For every fixed $n\ge1$, the maps
\begin{equation}\label{eq:return-Lone}
 R\longmapsto\widetilde J_{n,R}\in L^1(\nu;\R^4),\qquad
 R\longmapsto p_{n,R}\in L^1(\mathfrak m)
\end{equation}
are continuous. Both maps are uniformly continuous on $I$.
Moreover $\{\widetilde N_{n,R}:R\in I\}$ and
$\{\widetilde T_{n,R}:R\in I\}$ are uniformly integrable.
If $\eta_{n,R}$ is the record law, then
\begin{equation}\label{eq:record-Wone}
 W_1(\eta_{n,R},\eta_{n,S})
 \le c_*^{-1}\|\widetilde J_{n,R}-\widetilde J_{n,S}\|_{L^1(\nu)},
\end{equation}
where $W_1$ uses the Euclidean distance on $\Z^3\times\R$.
In particular, for the original, unmodified characteristic functions,
\begin{align}\label{eq:raw-frequency-comparison}
 |\Phi_{n,R}(\omega)-\Phi_{n,S}(\omega)|
 &\le \|p_{n,R}-p_{n,S}\|_{L^1(\mathfrak m)},\nonumber\\
 |\Phi_{n,R}(\omega)-\Phi_{n,S}(\omega)|
 &\le c_*^{-1}|\omega|\,
       \|\widetilde J_{n,R}-\widetilde J_{n,S}\|_{L^1(\nu)}.
\end{align}
No modulus uniform in $n$, and no frequency-integrated estimate, is
asserted by these inequalities.
\end{theorem}
\begin{proof}
The Kac identity and invariance of $F_R^*$ give, exactly,
\begin{equation}\label{eq:block-Kac-mass}
 \int_M\widetilde N_{n,R}\dd\nu=n,\qquad
 \int_M\widetilde T_{n,R}\dd\nu=n\bar\tau_R.
\end{equation}
For nonnegative functions $h_j\to h$ almost everywhere with
$\int h_j\to\int h$, the identity
\[
 \|h_j-h\|_1=\int h_j+\int h-2\int\min(h_j,h)
\]
and dominated convergence for $\min(h_j,h)\le h$ prove $L^1$
convergence. Apply this to the two nonnegative coordinates, using
Lemma~\ref{lem:finite-return-stability}, \eqref{eq:block-Kac-mass}
and continuity of $\bar\tau_R$. Uniform finite horizon bounds every
single-flight lattice label by a fixed constant. Thus
$|\widetilde K_{n,R}|\le C\widetilde N_{n,R}$. The already established
$L^1$ convergence of the counts makes these bounds uniformly integrable
along each convergent parameter sequence. Almost-everywhere convergence
and truncation then give $L^1$ convergence of the displacement as well.
This proves the first continuity in \eqref{eq:return-Lone}.

A continuous image of the compact interval $I$ is compact in $L^1$.
For completeness, an $L^1$-compact nonnegative family is uniformly
integrable. Cover it by finitely many $L^1$ balls of radius $\epsilon$
with centers $h_1,\ldots,h_q$. For a member $h$ in the ball about $h_i$,
\[
 \int_{\{h>A\}}h
 \le 2\|h-h_i\|_1+
              2\int_{\{h_i>A/2\}}h_i.
\]
Choose $A$ for the finite set of centers and then let $\epsilon$
decrease. This proves the stated uniform integrability without a tail
assumption or an exponential return estimate.

We next prove existence and $L^1$ continuity of the density. Partition
$Y_R^*$, up to null sets, by total physical count and ordered lattice
word. The argument of Corollary~\ref{cor:raw-density} applies to this
actual first-return event as well: on each regular word the cumulative
length has at most one critical point by
Proposition~\ref{prop:general-edge}; elsewhere it is a submersion.
Coarea on a countable exhaustion and then summation give absolute
continuity. This argument recomputes the return event for $Y_R^*$;
it does not transfer a density normalization from $Y_R$.

Fix $R_0$ and $\epsilon>0$. After deleting the null sets just described,
choose finitely many smooth nonnegative cutoffs $\psi_i$ on $M$, with
$\sum_i\psi_i\le1$, whose supports lie compactly in regular
first-$n$-return patches for $R_0$, such that
\[
 c_*^{-1}\int\sum_i\psi_i\dd\nu>1-\epsilon.
\]
Refine the patches so that on each one a coordinate derivative of the
cumulative length is nonzero. Compact support and the finite-return
lemma preserve the same word, all return decisions and this derivative
for $R$ close to $R_0$. On each patch the map consisting of the other
initial coordinate and the cumulative length is a local diffeomorphism.
It and its inverse depend continuously in $C^1$ on $R$ on slightly
larger compact patches. Changing variables and integrating the other
coordinate therefore shows that the pushforward of
$c_*^{-1}\psi_i\nu$ has a density continuous in $L^1$ with $R$.
One can see this directly by extending the transformed smooth weight
by zero off its coordinate patch: the compactly supported cutoff
vanishes near the patch boundary, the Jacobian is bounded away from
zero, and dominated convergence applies on a fixed bounded coordinate
box. The lattice coordinates are constant on the patch.

The remainder is a positive measure. Its mass is
$1-c_*^{-1}\int\sum_i\psi_i\dd\nu<\epsilon$, for all these $R$,
because the cutoff supports stay inside $Y_R^*$ and $\nu(Y_R^*)=c_*$.
Consequently
\[
 \|p_{n,R}-p_{n,R_0}\|_1
 \le\sum_i\|p_{i,R}-p_{i,R_0}\|_1+2\epsilon.
\]
Let $R\to R_0$ and then $\epsilon\to0$. This proves the second
continuity in \eqref{eq:return-Lone}. Compactness of $I$ gives uniform
continuity of both maps.

For a 1-Lipschitz function $f$, subtract the constant $f(0)$. Then
\[
 \int f\dd\eta_{n,R}
 =c_*^{-1}\int_M f(\widetilde J_{n,R})\dd\nu.
\]
The same identity holds for $S$. Subtraction and the Lipschitz bound,
followed by the dual characterization of $W_1$, give
\eqref{eq:record-Wone}. All first moments are finite by
\eqref{eq:block-Kac-mass}. Direct integration against the densities
gives the first bound in \eqref{eq:raw-frequency-comparison}. For the
second, use the zero extensions and
$|e^{ia}-e^{ib}|\le|a-b|$; the off-section constants cancel because
both section measures equal $c_*$. These are estimates of the raw
transforms, with no convolution or change of observable.
\end{proof}

\begin{proposition}[Weighted records and a positive denominator]
\label{prop:weighted-return-stability}
Fix $n$. Suppose measurable insertions $w_R:M\to\R$ satisfy
$|w_R|\le M_0$ and, for every $R_0$, $w_R(x)\to w_{R_0}(x)$ for
$\nu$-almost every $x$ as $R\to R_0$. The $w_R$-weighted record
densities $p^w_{n,R}$ are continuous in $L^1(\mathfrak m)$.
Set
\[
 d_n(\delta)=\sup_{|R-S|\le\delta}\|p_{n,R}-p_{n,S}\|_1,
 \qquad
 d_n^w(\delta)=\sup_{|R-S|\le\delta}\|p^w_{n,R}-p^w_{n,S}\|_1.
\]
These moduli tend to zero. For any Borel observation event $B$ in the
record space, if $\eta_{n,R}(B)\ge b>0$ and
$d_n(\delta)<b$, then, whenever $|R-S|\le\delta$,
\begin{equation}\label{eq:finite-conditioning-stability}
 \left|\frac{\int_Bp^w_{n,R}\dd\mathfrak m}{\eta_{n,R}(B)}
       -\frac{\int_Bp^w_{n,S}\dd\mathfrak m}{\eta_{n,S}(B)}\right|
 \le\frac{d_n^w(\delta)+M_0d_n(\delta)}{b-d_n(\delta)}.
\end{equation}
\end{proposition}
\begin{proof}
Use the compact regular patches in the preceding proof. Replacing
$w_R$ by $w_{R_0}$ costs at most
$c_*^{-1}\int|w_R-w_{R_0}|\dd\nu\to0$, by dominated convergence
and contraction of total variation under pushforward. On each fixed
patch approximate the bounded measurable weight $w_{R_0}$ in $L^1$
by a smooth one. The approximation error in its pushforward is at
most the same $L^1$ error, uniformly in $R$. For the smooth weight,
the coordinate-change proof above applies. The leftover measure has
total variation at most $M_0\epsilon$. This proves weighted density
continuity, and compactness gives the modulus.

Write $a_R=\int_Bp^w_{n,R}\dd\mathfrak m$ and
$b_R=\eta_{n,R}(B)$. Then
$|a_R-a_S|\le d_n^w(\delta)$,
$|b_R-b_S|\le d_n(\delta)$ and $|a_R|\le M_0b_R$.
Since $b_S\ge b-d_n(\delta)>0$, subtraction of the two ratios gives
\eqref{eq:finite-conditioning-stability}.
\end{proof}

The event $B$ may fix the lattice coordinates and restrict the roof
to an interval. The denominator in
\eqref{eq:finite-conditioning-stability} is an actual event probability.
The estimate does not remove its deterioration for small events. In
particular density $L^1$ continuity alone is not the shrinking-interval
LLT of Proposition~\ref{prop:conditioning}; that proposition retains
its stronger pointwise density hypotheses. An error bound
$|\widehat R-R|\le\delta$ can be inserted into the displayed moduli,
but no numerical sampling budget follows from a noneffective modulus.

\begin{corollary}[Uniform stationary endpoint control]
\label{cor:uniform-stationary-endpoints}
Put $H_R=\widetilde T_{1,R}$ and
$b_*=\min_{R\in I}\bar\tau_R>0$. The stationary induced suspensions,
regarded as probabilities on $M\times\R_+$, have densities
\[
 h_R(x,s)=\bar\tau_R^{-1}\one_{\{0\le s<H_R(x)\}}
\]
with respect to $\nu(\dd x)\dd s$, and
\begin{equation}\label{eq:suspension-Lone}
 \|h_R-h_S\|_1\le
 b_*^{-1}\bigl(\|H_R-H_S\|_1+|\bar\tau_R-\bar\tau_S|\bigr).
\end{equation}
Furthermore
\begin{equation}\label{eq:uniform-roof-tail}
 \Psi(q):=b_*^{-1}\sup_{R\in I}
                   \int_M H_R\one_{\{H_R>q\}}\dd\nu
             \longrightarrow0\quad(q\to\infty).
\end{equation}
Under the stationary flow, at either deterministic endpoint of an
interval of length $t$, the incomplete induced time, physical collision
count and lattice displacement are $o_P(\sqrt t)$, uniformly over
$R\in I$. The same is true jointly at any fixed finite number of
deterministic endpoints.
\end{corollary}
\begin{proof}
The normalization in \eqref{eq:induced-suspension} and
$c_*\int\tau_R^*\dd\nu_R^*=\bar\tau_R$ give $h_R$ exactly.
For fixed $x$, the symmetric difference of the intervals
$[0,H_R(x))$ and $[0,H_S(x))$ has length $|H_R(x)-H_S(x)|$.
Integrating, and then comparing the two normalizing constants, proves
\eqref{eq:suspension-Lone}. The uniform integrability in
Theorem~\ref{thm:return-stability} proves
\eqref{eq:uniform-roof-tail}.

The full roof containing a stationary endpoint has tail
\[
 \Pr_R\{\text{containing roof}>q\}
 =\bar\tau_R^{-1}\int_MH_R\one_{\{H_R>q\}}\dd\nu
 \le\Psi(q).
\]
This is the length-biased law; replacing it by
$\nu_R^*(\tau_R^*>q)$ would be incorrect. Both endpoint pieces
are at most this full roof. The collision count of an incomplete
piece of length $s$ is at most $1+(50/3)s$, and its lattice
displacement is bounded by $C(1+s)$, by the physical estimates in
Proposition~\ref{prop:induced-suspension}. Hence there are fixed
constants $C_0,C_1$ such that each marked endpoint norm is bounded
by $C_0+C_1H$ when its containing roof has length $H$. Its probability of exceeding
$\epsilon\sqrt t$ is at most
$\Psi((\epsilon\sqrt t-C_0)/C_1)$ for large $t$, uniformly in $R$.
Stationarity applies at every deterministic endpoint, including an
endpoint depending deterministically on $t$. A union bound proves
the finite-endpoint assertion; independence is not used.
\end{proof}

This removes the additional uniform-integrability premise for fixed
stationary endpoints in this family. It is not a bound on a supremum
over a growing time interval or at an arbitrary stopping time. The
functional clock theorem still needs its printed process hypotheses.
Likewise \eqref{eq:raw-frequency-comparison} is a uniform comparison
of raw characteristic functions, not an integrable tail for any one
of them. The high-frequency and central all-branch remainder problem
is unchanged by an $L^1$ comparison of probability densities.
